\chapter{Pendahuluan}\label{chap:metr}


\section{Apakah pendekatan aksiomatik itu?}
\label{preaxioms}

Dalam pendekatan aksiomatik, bidang didefinisikan sebagai apa pun yang memenuhi seperangkat sifat tertentu yang disebut {}\emph{aksioma}.
Aksioma-aksioma ini pada dasarnya berperan sebagai aturan permainan matematika.
Setelah sistem aksioma ditetapkan, suatu pernyataan benar jika pernyataan itu mengikuti secara logis dari aksioma-aksioma tersebut, dan hanya dalam hal itu.

Perumusan aksioma-aksioma pertama sama sekali tidak ketat.
Sebagai contoh, Euklides mendeskripsikan sebuah {}\emph{garis} sebagai \textit{panjang tanpa lebar}
dan sebuah {}\emph{garis lurus} sebagai garis yang \textit{terletak merata dengan titik-titik pada dirinya}.
Namun, perumusan ini cukup jelas bagi para matematikawan untuk saling memahami.

Cara terbaik untuk memahami suatu sistem aksiomatik
adalah dengan menyusunnya sendiri.
Untuk saat ini, anggaplah kita memahami secara intuitif konsep-konsep seperti garis dan titik.
Bayangkan suatu permukaan papan tulis yang tak berhingga dan sempurna. 
Marilah kita mencoba menghimpun pengamatan-pengamatan utama mengenai model ini.

\begin{enumerate}[(i)]
 \item\label{preaxiomI} Kita dapat mengukur jarak antartitik.
 \item\label{preaxiomII} Kita dapat menggambar tepat satu garis yang melalui dua titik yang diberikan.
 \item\label{preaxiomIII} Kita dapat mengukur sudut.
 \item\label{preaxiomIV} Jika kita memutar atau menggeser, kita tidak akan melihat perbedaan.
 \item\label{preaxiomV} Jika kita mengubah skala, kita tidak akan melihat perbedaan.
\end{enumerate}
Pengamatan-pengamatan ini merupakan awal yang baik.
Kelak, kita akan mengembangkan bahasa untuk merumuskannya kembali secara ketat.

\section{Apakah model itu?}
\label{page:model}

Bidang Euklides dapat didefinisikan secara ketat sebagai berikut:

\textit{Definisikan sebuah {}\emph{titik} pada bidang Euklides sebagai pasangan bilangan real $(x,y)$ dan definisikan {}\emph{jarak} antara dua titik $(x_1,y_1)$ dan $(x_2,y_2)$ dengan rumus berikut:}
\[\sqrt{(x_1-x_2)^2+(y_1-y_2)^2}.\]

Itu saja!
Kita telah memberikan sebuah {}\emph{model numerik} untuk bidang Euklides;
model tersebut membangun bidang Euklides dari bilangan real, yang dianggap telah dikenal.

Keringkasan merupakan keunggulan utama pendekatan model,
tetapi secara intuitif tidak jelas mengapa kita mendefinisikan titik dan jarak dengan cara ini.

Di sisi lain, pengamatan-pengamatan pada bagian sebelumnya jelas secara intuitif ---
inilah keunggulan utama pendekatan aksiomatik.
Keunggulan lainnya --- pendekatan aksiomatik mudah disesuaikan. 
Sebagai contoh, kita dapat menghapus atau mengganti sebuah aksioma dari daftar. 
Kita akan melakukannya di Bab~\ref{chap:non-euclid}.

\section{Ruang metrik}\label{sec:Metric spaces}

Konsep ruang metrik memberikan 
cara yang ketat untuk mengatakan: \textit{``kita dapat mengukur jarak antartitik''}.
Dengan kata lain, sebagai pengganti (\ref{preaxiomI}) pada Bagian~\ref{preaxioms},
kita dapat mengatakan \textit{``bidang Euklides adalah ruang metrik''}.

\begin{thm}{Definisi}\label{def:metric-space}
Misalkan $\mathcal X$ suatu himpunan tak kosong dan 
misalkan $d$ suatu fungsi
yang menetapkan bilangan real $d(A,B)$
untuk setiap pasangan $A,B\in\mathcal X$.
Maka $d$ disebut \index{metrik}\emph{metrik} pada 
$\mathcal X$ jika, untuk setiap
$A,B,C\in \mathcal X$, syarat-syarat berikut terpenuhi:
\begin{enumerate}[(a)]
\item\label{def:metric-space:a} Kepositifan: 
$d(A,B)\ge 0.$
\item\label{def:metric-space:b} $A=B$ jika dan hanya jika 
$d(A,B)=0.$
\item\label{def:metric-space:c} Simetri: $d(A, B) = d(B, A).$
\item\label{def:metric-space:d} \index{segitiga!ketaksamaan}Ketaksamaan segitiga: 
$d(A, C) \le d(A, B) + d(B, C).$
\end{enumerate}
Sebuah \index{metrik!ruang}\emph{ruang metrik} didefinisikan sebagai himpunan yang dilengkapi dengan metrik. 
Secara lebih formal, ruang metrik adalah pasangan $(\mathcal X, d)$ dengan $\mathcal X$ suatu himpunan dan $d$ suatu metrik pada~$\mathcal X$.

Unsur-unsur $\mathcal X$ disebut \index{titik}\emph{titik-titik} ruang metrik tersebut.
Nilai $d(A, B)$ disebut \index{jarak}\emph{jarak} dari $A$ ke~$B$.
\end{thm}

\pagebreak

\subsection*{Contoh-contoh}

\begin{itemize}
\item {}\emph{Metrik diskret.} Misalkan $\mathcal X$ suatu himpunan sembarang. 
Untuk setiap $A,B\z\in\mathcal X$, tetapkan $d(A,B)\z=0$ jika $A=B$ dan $d(A,B)=1$ jika tidak.
Metrik $d$ disebut \index{metrik diskret}\emph{metrik diskret} pada~$\mathcal X$.
\item\index{real!garis bilangan}\emph{Garis bilangan real.}
Himpunan semua bilangan real ($\mathbb{R}$) dengan metrik $d$ yang didefinisikan sebagai 
$$d(A,B)\df|A-B|.$$
\end{itemize}

\begin{thm}{Latihan}\label{ex:dist-square}
Tunjukkan bahwa $d(A,B)=|A-B|^2$ \textit{bukan} metrik pada $\mathbb{R}$.
\end{thm}

\begin{itemize}
\item \textit{Metrik pada bidang.}
Misalkan $\mathbb{R}^2$ menyatakan himpunan semua pasangan $(x,y)$ bilangan real.
Andaikan $A=(x_A,y_A)$ dan $B=(x_B,y_B)$.
Perhatikan metrik-metrik berikut pada $\mathbb{R}^2$:
\begin{itemize}
\item\index{Euklides!metrik}\emph{Metrik Euklides,} dilambangkan dengan \index{59@$d_1$, $d_2$, $d_\infty$}$d_2$, dan didefinisikan sebagai \label{def:d_2}
$$d_2(A,B)=\sqrt{(x_A-x_B)^2+(y_A-y_B)^2}.$$
\item\label{Manhattan plane}\index{bidang Manhattan}\emph{Metrik Manhattan,} dilambangkan dengan $d_1$ dan didefinisikan sebagai 
$$d_1(A,B)=|x_A-x_B|+|y_A-y_B|.$$
\item{}\emph{Metrik maksimum,} dilambangkan dengan $d_\infty$ dan didefinisikan sebagai 
$$d_\infty(A,B)=\max\{|x_A-x_B|,|y_A-y_B|\}.$$
\end{itemize}
\end{itemize}

\begin{thm}{Latihan}\label{ex:d_1+d_2+d_infty}
Buktikan bahwa fungsi-fungsi berikut merupakan metrik pada $\mathbb{R}^2$:
(a)~$d_1$; (b)~$d_2$; (c)~$d_\infty$.
\end{thm}


\section{Notasi singkat untuk jarak}

Biasanya, kita hanya memerlukan satu metrik pada suatu ruang.
Oleh karena itu, kita tidak perlu menyebutkan metrik tersebut setiap kali.

Diberikan suatu ruang metrik $\mathcal X$,
jarak antara titik $A$ dan $B$ akan dilambangkan dengan
$$AB
\quad
\text{atau}
\quad
d_{\mathcal X}(A,B);$$
notasi yang terakhir hanya digunakan jika kita perlu menegaskan bahwa $A$ dan $B$ merupakan titik-titik ruang metrik~$\mathcal X$.

Sebagai contoh, ketaksamaan segitiga dapat ditulis sebagai 
$$AC\le AB+BC.$$

Untuk menghindari kerancuan, kita akan selalu menggunakan ``$\cdot$'' untuk perkalian,
agar $AB$ tidak disalahartikan sebagai $A\cdot B$.

\begin{thm}{Latihan}\label{ex:4-triangle}
Tunjukkan bahwa ketaksamaan
\[AB+PQ\le AP+AQ+BP+BQ\]
berlaku untuk sembarang empat titik $A$, $B$, $P$, $Q$ dalam suatu ruang metrik.
\end{thm}


\section{Isometri, transformasi isometrik, dan garis}

Pada bagian ini, kita mendefinisikan garis dalam ruang metrik.
Setelah itu dilakukan, kalimat \textit{``Kita dapat menggambar tepat satu garis yang melalui dua titik yang diberikan.''} menjadi ketat; lihat (\ref{preaxiomII}) pada Bagian~\ref{preaxioms}. 

Ingatlah bahwa suatu pemetaan $f\:\mathcal{X}\to\mathcal{Y}$
merupakan \index{bijeksi}\emph{bijeksi}
jika pemetaan tersebut memasangkan secara tepat unsur-unsur kedua himpunan.
Secara ekuivalen, $f\:\mathcal{X}\to\mathcal{Y}$ merupakan bijeksi jika memiliki suatu \index{invers}\emph{invers};
yaitu, suatu pemetaan $g\:\mathcal{Y}\to\mathcal{X}$
sedemikian sehingga 
$g(f(A))\z=A$ untuk setiap $A\in\mathcal{X}$
dan
$f(g(B))\z=B$ untuk setiap $B\in\mathcal{Y}$. 

Misalkan $\mathcal X$ dan $\mathcal Y$ dua ruang metrik serta $d_{\mathcal X}$ dan $d_{\mathcal Y}$ metriknya. 
Suatu pemetaan 
$$f\:\mathcal X \z\to \mathcal Y$$ 
disebut
\index{pemetaan yang mempertahankan jarak}\emph{pemetaan yang mempertahankan jarak} jika 
$$d_{\mathcal Y}(f(A), f(B))
 = d_{\mathcal X}(A,B)$$
untuk semua $A,B\in {\mathcal X}$.

Pemetaan bijektif yang mempertahankan jarak disebut \index{isometri}\emph{isometri}. 

Dua ruang metrik disebut \index{ruang isometrik}\emph{isometrik} jika terdapat suatu isometri di antara keduanya.

Suatu isometri dari ruang metrik ke dirinya sendiri juga disebut \index{transformasi isometrik}\emph{transformasi isometrik}.

\begin{thm}{Latihan}\label{ex:dist-preserv=>injective}
Tunjukkan bahwa setiap pemetaan yang mempertahankan jarak bersifat \index{pemetaan injektif}\emph{injektif};
yaitu, jika $f\:\mathcal X\to\mathcal Y$ merupakan pemetaan yang mempertahankan jarak, 
maka $f(A)\ne f(B)$ untuk setiap pasangan titik berbeda $A, B\in \mathcal X$.
\end{thm}

\begin{thm}[!]{Latihan}\label{ex:motion-of-R}
Tunjukkan bahwa jika $f\:\mathbb{R}\to\mathbb{R}$ merupakan transformasi isometrik pada garis bilangan real,
maka salah satu dari (a)
$f(x)=f(0)+x$ untuk setiap $x\in \mathbb{R}$, 
atau (b)
$f(x)=f(0)-x$ untuk setiap $x\in \mathbb{R}$. 

\end{thm}

\begin{thm}{Latihan}\label{ex:d_1=d_infty}
Buktikan bahwa $(\mathbb{R}^2,d_1)$ isometrik dengan $(\mathbb{R}^2,d_\infty)$.
\end{thm}

\begin{thm}{Latihan lanjutan}\label{ad-ex:motions of Manhattan plane}
Uraikan semua transformasi isometrik bidang Manhattan (didefinisikan dalam Bagian~\ref{sec:Metric spaces}).
\end{thm}

Jika $\mathcal X$ ruang metrik dan $\mathcal Y$ subhimpunan dari $\mathcal X$,
maka suatu metrik pada $\mathcal Y$ dapat diperoleh dengan membatasi metrik dari~$\mathcal X$. 
Dengan kata lain, 
jarak antara dua titik di $\mathcal Y$ didefinisikan sebagai jarak antara kedua titik tersebut di $\mathcal X$.
Dengan cara ini, setiap subhimpunan ruang metrik juga dapat dipandang sebagai ruang metrik.

\begin{thm}{Definisi}\label{def:line}
Suatu subhimpunan $\ell$ dari ruang metrik disebut \index{garis}\emph{garis} jika isometrik dengan garis bilangan real.
\end{thm}

Tripel titik yang terletak pada satu garis disebut \index{titik-titik segaris}\emph{segaris}.
Perhatikan bahwa jika $A$, $B$, dan  $C$ segaris, $AC\ge AB$, dan $AC\ge BC$, maka $AC\z= AB+BC$.

Beberapa ruang metrik tidak memiliki garis; sebagai contoh, ruang dengan metrik diskret.
Gambar menunjukkan contoh-contoh garis pada bidang Manhattan $(\mathbb{R}^2,d_1)$. 
\begin{figure}[!ht]
\centering
\includegraphics{mppics/pic-2}
\end{figure}

\begin{thm}{Latihan}\label{ex:y=|x|}
Perhatikan grafik $y=|x|$ dalam $\mathbb{R}^2$.
Dalam ruang-ruang berikut ini,
(a) $(\mathbb{R}^2,d_1)$, 
(b) $(\mathbb{R}^2,d_2)$, 
(c) $(\mathbb{R}^2,d_\infty)$ 
apakah grafik tersebut membentuk sebuah garis? 
Mengapa?
\end{thm}

\begin{thm}[!]{Latihan}\label{ex:line-motion}
Tunjukkan bahwa setiap transformasi isometrik memetakan sebuah garis ke sebuah garis.
\end{thm}


\section{Sinar dan ruas}

Andaikan terdapat sebuah garis $\ell$ yang melalui
dua titik berbeda $P$ dan $Q$.
Dalam hal ini, kita dapat menyatakan $\ell$ sebagai $(PQ)$.
Mungkin terdapat lebih dari satu garis yang melalui $P$ dan $Q$,
tetapi jika kita menulis \index{60@$(PQ)$, $[PQ)$, $[PQ]$}$(PQ)$, kita menganggap bahwa kita telah memilih salah satu garis tersebut. 

Kita akan menggunakan $[P Q)$ untuk menyatakan \index{sinar}\emph{sinar}
yang berawal di $P$ dan memuat~$Q$. 
Secara formal, $[P Q)$ merupakan subhimpunan dari $(P Q)$ yang bersesuaian dengan $[0,\infty)$ melalui suatu isometri $f\:(P Q)\to \mathbb{R}$ sedemikian sehingga $f(P)=0$ dan $f(Q)>0$.

Subhimpunan garis $(P Q)$ yang berada di antara $P$ dan $Q$ disebut \index{ruas}\emph{ruas antara} $P$ dan $Q$; ruas tersebut dilambangkan dengan~$[P Q]$.
Secara formal, ruas dapat didefinisikan sebagai irisan dua sinar: $[P Q]=[P Q)\cap[Q P)$.

\begin{thm}[!]{Latihan}\label{ex:trig==}
Tunjukkan bahwa, untuk $r\ge 0$ dan suatu sinar $[P Q)$ yang diberikan, terdapat tepat satu $X\in [P Q)$  sedemikian sehingga $P X=r$.
%\begin{enumerate}[(a)]
%\item if $X\in [PQ)$, then $QX=|PX-PQ|$;
%\item if $X\in [PQ]$, then $PX+XQ=PQ$;
%\item $[OP)\cup[OQ)$ is a line if and only if $XO+OY=XY$ for all $X\in [OP)$ and $Y\in[OQ)$.
%\end{enumerate}

\end{thm}

%\begin{thm}{Exercise}\label{ex:lineABC}
%Let $A$, $B$, and $C$ be disticnt points in a metric space such that $AB+BC\z=AC$.
%Suppose that the there are lines $(AB)$, $(AC)$, and $(BC)$.
%Show that $((AC)\setminus[AC])\cup [AB]\cup [BC]$ is a line.
%\end{thm}

\section{Sudut}

Tujuan kita selanjutnya adalah memperkenalkan sudut dan ukuran sudut; 
setelah itu, pernyataan \textit{``kita dapat mengukur sudut''} akan menjadi ketat;
lihat (\ref{preaxiomIII}) pada Bagian~\ref{preaxioms}.

Pasangan terurut dua sinar yang berawal di titik yang sama disebut \index{sudut}\emph{sudut}.
Sudut $AOB$ (juga dilambangkan dengan \index{10@$\angle$, $\measuredangle$}$\angle AOB$) adalah pasangan sinar $[OA)$ dan $[OB)$; titik $O$ disebut \index{titik sudut!sudut}\emph{titik sudut} dari sudut tersebut.

Secara intuitif, ukuran sudut menyatakan seberapa jauh sinar pertama harus diputar berlawanan arah jarum jam,
agar mencapai posisi sinar kedua.
Satu putaran penuh dianggap sebesar $2\cdot\pi$;
nilai ini bersesuaian dengan ukuran sudut dalam radian.%
\footnote{Untuk saat ini, kita dapat memandang $\pi$ sebagai
ukuran setengah putaran dalam satuan tertentu.
Nilai numeriknya, $\pi\approx 3.14$, tidak penting pada tahap ini.}

Ukuran sudut $\angle AOB$ dilambangkan dengan $\measuredangle AOB$;
nilainya merupakan bilangan real dalam interval $(-\pi,\pi]$. 

\begin{wrapfigure}{o}{25mm}
\vskip-0mm
\centering
\includegraphics{mppics/pic-4}
\end{wrapfigure}

Notasi $\angle AOB$ dan $\measuredangle AOB$ tampak serupa;
keduanya juga memiliki makna yang berkaitan erat tetapi berbeda dan sebaiknya tidak dirancukan.
Sebagai contoh, persamaan 
$\angle AOB\z=\angle A'O'B'$
berarti bahwa
$[OA)=[O'A')$ dan $[OB)\z=[O'B')$;
khususnya, $O=O'$.
Di sisi lain, persamaan 
$\measuredangle AOB\z=\measuredangle A'O'B'$ 
hanya menyatakan kesamaan dua bilangan real;
dalam hal ini, $O$ mungkin berbeda dari~$O'$.

Berikut adalah sifat pertama ukuran sudut yang akan menjadi bagian dari aksioma.
\textit{Diberikan suatu sinar $[O A)$ dan $\alpha\in(-\pi,\pi]$, terdapat tepat satu sinar $[O B)$ sedemikian sehingga $\measuredangle A O B= \alpha$.}



\section[\texorpdfstring{Bilangan real modulo $2\cdot\pi$}{Bilangan real modulo 2·π}]{Bilangan real modulo $\bm{2\cdot\pi}$}



Perhatikan tiga sinar yang berawal di titik yang sama, $[O A)$, $[O B)$, dan $[O C)$.
Ketiganya membentuk tiga sudut $A O B$, $B O C$, dan $A O C$,
sehingga nilai $\measuredangle A O C$ seharusnya sama dengan
jumlah $\measuredangle A O B+\measuredangle B O C$ modulo satu putaran penuh.
Sifat ini akan dinyatakan dengan rumus 
$$\measuredangle A O B+\measuredangle B O C\equiv \measuredangle A O C,$$
di mana \index{34@$\equiv$}``$\equiv$'' adalah notasi baru yang akan segera kita perkenalkan.
Identitas terakhir ini akan menjadi bagian dari aksioma-aksioma.

Kita akan menulis $\alpha\equiv\beta\pmod{2\cdot\pi}$, atau secara ringkas
\begin{align*}
\alpha&\equiv\beta
\end{align*}
jika $\alpha=\beta+2\cdot\pi\cdot n$
untuk suatu bilangan bulat~$n$.
Dalam hal ini, kita mengatakan 
$$\textit{``$\alpha$ sama dengan $\beta$ modulo $2\cdot\pi$''}.$$
Sebagai contoh, $-\pi
\equiv
\pi\equiv 3\cdot\pi$ dan $\tfrac12\cdot\pi
\equiv
-\tfrac32\cdot\pi$.

Relasi ``$\equiv$'' yang diperkenalkan berperilaku seperti tanda sama dengan,
tetapi
\[\dots\equiv\alpha-2\cdot\pi\equiv \alpha\equiv \alpha+2\cdot\pi\equiv \alpha+4\cdot\pi\equiv\dots;\] 
yaitu, jika dua ukuran sudut berbeda sebesar kelipatan satu putaran penuh,
maka keduanya dianggap sama.

Dengan ``$\equiv$'', kita dapat melakukan penjumlahan, pengurangan, dan perkalian dengan bilangan bulat tanpa menimbulkan masalah.
Artinya, jika
$$\alpha\equiv\beta
\quad
\text{dan}
\quad
\alpha'\equiv \beta',$$ 
maka
$$\alpha+\alpha'\equiv\beta+\beta',
\quad
\alpha-\alpha'\equiv \beta-\beta'
\quad 
\text{dan}
\quad
n\cdot\alpha\equiv n\cdot\beta$$
untuk setiap bilangan bulat~$n$.
Namun, ``$\equiv$'' tidak mempertahankan perkalian dengan bilangan nonbulat; sebagai contoh, 
$$\pi
\equiv 
-\pi
\quad
\text{tetapi}
\quad
\tfrac12\cdot\pi
\not\equiv
-\tfrac12\cdot\pi.$$ 

\begin{thm}[!]{Latihan}\label{ex:2a=0}
Tunjukkan bahwa $2\cdot\alpha\equiv0$ jika dan hanya jika $\alpha\equiv0$ atau $\alpha\equiv\pi$.
\end{thm}

\begin{thm}{Latihan}\label{ex:a+b==c}
Andaikan $0\le \alpha\le \pi$, $0\le \beta\le \pi$, dan $0<\gamma<2\cdot \pi$.
Tunjukkan bahwa
\[\alpha+\beta\equiv \gamma\quad\text{mengakibatkan}\quad\alpha+\beta=\gamma.\]
\end{thm}


\section{Kekontinuan}

Ukuran sudut juga diasumsikan kontinu.
Tepatnya, sifat ukuran sudut berikut akan menjadi bagian dari aksioma-aksioma:

\textit{Fungsi}
$$\measuredangle\:(A,O,B)\mapsto\measuredangle A O B$$
\textit{kontinu pada setiap tripel titik $(A,O,B)$
sedemikian sehingga $O\ne A$, $O\ne B$, dan $\measuredangle A O B\ne\pi$.}

Untuk memperjelas sifat ini, kita perlu memperluas konsep {}\emph{kekontinuan} ke fungsi-fungsi antara ruang-ruang metrik.
Definisi ini merupakan perumuman langsung dari definisi baku untuk fungsi real ke real.

Misalkan $\mathcal X$ dan $\mathcal Y$ dua ruang metrik,
serta $d_{\mathcal X}$ dan $d_{\mathcal Y}$ metriknya.

Suatu pemetaan $f\:\mathcal X\to\mathcal Y$ disebut \index{kontinu}\emph{kontinu} di titik $A\in \mathcal X$
jika, untuk setiap $\epsilon>0$, terdapat $\delta>0$ sedemikian sehingga 
\[d_{\mathcal X}(A,A')
<
\delta
\quad
\Rightarrow
\quad
d_{\mathcal Y}(f(A),f(A'))
<
\epsilon.\]
(Secara informal, ini berarti bahwa perubahan $A$ yang cukup kecil menghasilkan perubahan $f(A)$ yang sekecil apa pun.)

Suatu pemetaan $f\:\mathcal X\to\mathcal Y$ disebut \index{kontinu}\emph{kontinu} jika kontinu di setiap titik $A\in \mathcal X$.

Pemetaan kontinu dengan beberapa variabel dapat didefinisikan dengan cara serupa.
Andaikan $f$ suatu fungsi yang menghasilkan sebuah titik di $\mathcal Y$ dari suatu tripel titik $(A,B,C)$
dalam~$\mathcal X$.
Pemetaan $f$ mungkin hanya didefinisikan untuk beberapa tripel dalam~$\mathcal X$.

Andaikan $f(A,B,C)$ terdefinisi.
Kemudian, kita mengatakan bahwa $f$ kontinu pada tripel $(A,B,C)$ 
jika, untuk setiap $\epsilon>0$, terdapat $\delta>0$ sedemikian sehingga
\[d_{\mathcal Y}(f(A,B,C),f(A',B',C'))<\epsilon\]
setiap kali $d_{\mathcal X}(A,A')<\delta$, $d_{\mathcal X}(B,B')<\delta$, dan $d_{\mathcal X}(C,C')<\delta$.


\begin{thm}[!]{Latihan}\label{ex:dist-cont}
Misalkan $\mathcal{X}$ suatu ruang metrik.
\begin{enumerate}[(a)]
\item\label{ex:dist-cont:a} Misalkan $A\in \mathcal{X}$ suatu titik tetap.
Tunjukkan bahwa fungsi 
$$f\:B\mapsto
d_{\mathcal{X}}(A,B)$$ 
kontinu di setiap titik~$B$.
\item Tunjukkan bahwa fungsi $g\:(A,B)\mapsto d_{\mathcal{X}}(A,B)$ kontinu pada setiap pasangan $A,B\in \mathcal{X}$.
\end{enumerate}

\end{thm}

\begin{thm}[!]{Latihan}\label{ex:comp+cont}
Misalkan $\mathcal{X}$, $\mathcal{Y}$, dan $\mathcal{Z}$ ruang-ruang metrik.
Andaikan fungsi $f\:\mathcal{X}\to\mathcal{Y}$
dan $g\:\mathcal{Y}\to\mathcal{Z}$ kontinu di setiap titik,
dan $h=g\circ f$ merupakan komposisi keduanya;
yaitu, $h(A)=g(f(A))$ untuk setiap $A\in \mathcal{X}$.
Tunjukkan bahwa $h\:\mathcal{X}\to\mathcal{Z}$ kontinu di setiap titik.
\end{thm}

\begin{thm}{Latihan}\label{ex:isom-cont}
Tunjukkan bahwa setiap pemetaan yang mempertahankan jarak bersifat kontinu.
\end{thm}

\section{Segitiga kongruen}
\label{sec:cong-triangles}

Tujuan kita selanjutnya adalah menemukan interpretasi yang ketat bagi pernyataan (\ref{preaxiomIV}) pada Bagian~\ref{preaxioms}.
Untuk itu, kita akan memperkenalkan konsep segitiga kongruen,
sehingga sebagai pengganti \textit{``jika kita memutar atau menggeser, kita tidak akan melihat perbedaan''}, kita dapat mengatakan bahwa kekongruenan sisi-sudut-sisi berlaku untuk segitiga.

Suatu tripel \textit{terurut} dari titik-titik \textit{berbeda} dalam ruang metrik $\mathcal{X}$, 
misalnya $A$, $B$, $C$,
disebut \index{segitiga}\emph{segitiga $ABC$}\label{page:def:triangle} (secara ringkas \index{20@$\triangle$}$\triangle A B C$).
Titik-titik $A$, $B$, dan $C$ disebut \index{titik sudut!segitiga}\emph{titik-titik sudut} dari $\triangle ABC$.
Perhatikan bahwa segitiga $A B C$ dan $A C B$ berbeda.

Dua segitiga $A' B' C'$ dan $A B C$ bersifat  
\index{segitiga!segitiga kongruen}
\index{kongruen!segitiga}\emph{kongruen}
(secara ringkas  \index{32@$\cong$}$\triangle A' B' C'\z\cong\triangle A B C$) jika terdapat transformasi isometrik $f\:\mathcal{X}\to\mathcal{X}$ sedemikian sehingga 
\[A'\z=f(A),
\quad
B'=f(B)
\quad
\text{dan}
\quad
C'=f(C).\]

Misalkan $\mathcal X$ suatu ruang metrik,
dan $f,g\:\mathcal X\to\mathcal X$ dua transformasi isometrik.
Perhatikan bahwa invers $f^{-1}:\mathcal X\to\mathcal X$,
serta komposisi $f\circ g:\mathcal X\to\mathcal X$,
juga merupakan transformasi isometrik.

Oleh karena itu, ``$\cong$'' merupakan \index{relasi ekuivalensi}\emph{relasi ekuivalensi};
yaitu, setiap segitiga kongruen dengan dirinya sendiri,
dan dua syarat berikut berlaku:
\begin{itemize} 
\item Jika $\triangle A' B' C'\z\cong\triangle A B C$, maka $\triangle A B C\z\cong\triangle A' B' C'$.
\item Jika $\triangle A'' B'' C''\z\cong\triangle A' B' C'$ dan $\triangle A' B' C'\z\cong\triangle A B C$,
maka 
$$\triangle A'' B'' C''\cong\triangle A B C.$$
\end{itemize}

Jika $\triangle A' B' C'\z\cong\triangle A B C$,
maka $AB\z=A'B'$,
$BC=B'C'$ dan $CA=C'A'$.

Untuk metrik diskret, serta beberapa metrik lainnya, 
kebalikannya juga berlaku.
Contoh berikut menunjukkan bahwa hal tersebut tidak berlaku pada bidang Manhattan:

\parbf{Contoh.}\label{example:isometric but not congruent}
Perhatikan tiga titik 
$A=(0,1)$, $B=(1,0)$, dan $C\z=(-1,0)$ pada bidang Manhattan $(\mathbb{R}^2,d_1)$.
Maka
$$d_1(A,B)=d_1(A,C)=d_1(B,C)=2.$$

Di satu pihak, $\triangle ABC\cong \triangle ACB$.
Memang, pemetaan $(x,y)\z\mapsto (-x,y)$ merupakan transformasi isometrik pada $(\mathbb{R}^2,d_1)$
yang memetakan $A\z\mapsto A$, $B\mapsto C$, dan $C\z\mapsto B$.

Di pihak lain, $\triangle ABC\ncong \triangle BCA$.
Memang, untuk memperoleh kontradiksi, andaikan $\triangle ABC\cong \triangle BCA$; yaitu, terdapat transformasi isometrik $f$ pada $(\mathbb{R}^2,d_1)$ yang memetakan $A\mapsto B$, $B\mapsto C$, dan $C\mapsto A$.

\begin{wrapfigure}[6]{o}{33mm}
\vskip-4mm
\centering
\includegraphics{mppics/pic-6}
\end{wrapfigure}

Kita mengatakan bahwa $M$ merupakan titik tengah $A$ dan $B$ jika 
\[d_1(A,M)=d_1(B,M)=\tfrac12\cdot d_1(A,B).\]
Perhatikan bahwa suatu titik $M$ merupakan titik tengah $A$ dan $B$ jika dan hanya jika $f(M)$ merupakan titik tengah $B$ dan~$C$.

Himpunan titik tengah $A$ dan $B$ tak berhingga; himpunan tersebut memuat semua titik $(t,t)$ untuk $t\in[0,1]$ (himpunan itu merupakan ruas berwarna abu-abu pada gambar di atas).
Sebaliknya, titik tengah $B$ dan $C$ tunggal.
Oleh karena itu, pemetaan $f$ tidak mungkin bijektif sehingga menghasilkan kontradiksi.

\begin{thm}{Latihan}\label{ex:ncong}
Perhatikan suatu ruang metrik yang terdiri atas empat titik $A$, $B$, $C$, $D$ dengan metrik yang didefinisikan oleh $AB=AC=AD=BC=BD=1$, dan 
$CD=2$.
Tunjukkan bahwa (a) $\triangle ABC\cong \triangle BAC$, dan (b) $\triangle ABC\ncong \triangle BCA$.
\end{thm}

